\section{A Reproducible Diagnostic Walkthrough}
\label{sec:case}
The benchmark diagnoses aggregate behavior; the following case demonstrates a user's complete inspection and intervention. We constructed two short documents about a fictional Harbor sensor pilot and executed the installed reviewer package with real generation and embedding APIs. This new deployment case is separate from the historical benchmark. The question is: ``For how many days does the Harbor sensor pilot retain raw sensor logs?'' The operating guide refers to a separate data policy; only that policy gives the answer, 17 days.

\paragraph{Trigger and inspection.}
With only the operating guide selected, the user receives the generic message ``MARA could not retrieve enough evidence to answer reliably.'' Opening the structured response, however, shows \texttt{route=doc\_text} and retrieval \texttt{status=good}. The captured provider call also completed. The blocking decision is instead verification \texttt{status=unknown}, with guardrail \texttt{action=abstain}. The evidence bundle names the operating guide and preserves its sentence directing the reader to the data policy (Figure~\ref{fig:diagnosis}). Thus the message alone does not accurately localize the failed stage: retrieving a relevant document did not establish support for the requested duration.

\begin{figure}[t]
\centering
\begin{tikzpicture}[font=\small,node distance=3mm,
box/.style={draw,rounded corners,align=left,text width=7.0cm,inner sep=4pt},
arr/.style={-{Stealth[length=2mm]}}]
\node[box] (a) {\textbf{1. Task unresolved}\\Operating guide selected; user receives a refusal.};
\node[box,below=of a] (b) {\textbf{2. Inspect the recorded stages}\\Text route; retrieval: \texttt{good}; provider call completed.\\Verification: \texttt{unknown}; guardrail: \texttt{abstain}.};
\node[box,below=of b] (c) {\textbf{3. Locate the missing support}\\Selected evidence: ``Data retention is specified in the separate Harbor data policy.''};
\node[box,below=of c] (d) {\textbf{4. Intervene}\\Select the data policy; repeat the same question.\\Keep the route policy, provider and strict verification.};
\node[box,below=of d] (e) {\textbf{5. Check the result}\\Answer: ``exactly 17 days'' with a citation.\\Retrieval: \texttt{good}; verification: \texttt{supported};\\guardrail: \texttt{return}. Text route remains selected.};
\draw[arr] (a)--(b); \draw[arr] (b)--(c);
\draw[arr] (c)--(d); \draw[arr] (d)--(e);
\end{tikzpicture}
\caption{A walkthrough derived from captured runtime records, not a new UI screenshot. Complete runtime responses and captured provider answer texts are supplied. The intervention changes source selection; it is not an automatic route switch.}
\label{fig:diagnosis}
\end{figure}

\paragraph{Action and observed outcome.}
The user selects the data policy and repeats the question in a fresh conversation. The serialized requests differ only in selected source IDs; effective settings are equal. The second response states ``exactly 17 days'' with a citation, verification becomes \texttt{supported}, and the guardrail returns the answer. Both turns use text retrieval. An independent repetition instead rejects extra generated explanation even with the correct source; both runs are retained (Appendix~\ref{app:case}). Thus this observed source-scope correction is not a reliable-recovery claim. The generic refusal also conflates lack of answer support with retrieval failure; we retain that discrepancy.

\paragraph{What the extension enables.}
Kotaemon already supports selecting another file and inspecting evidence. DocQA-Inspect's added value in this task is to distinguish the selected route, retrieval success, generation, verification, and final guardrail, then export these observations together. Table~\ref{tab:boundary} compares the inspection operations against pinned upstream source. It is a source-based walkthrough, not a timed usability study or a performance comparison; we make no claim of fewer clicks or faster diagnosis.

\paragraph{Auditable answers and configuration.}
The case package preserves both raw provider outputs, the answers before verification and guardrails, the displayed answers, complete requests and effective settings, selected evidence, and controller traces. An observer records the actual text-provider calls without altering their inputs or output chunks. The displayed and scoring answers are identical under this case's declared scoring rule, and the supplied script recomputes its checks from those strings. This closes the raw-output audit path for the new case only. It cannot recover the historical SlideVQA answers or establish why their scores tie. Appendix~\ref{app:case} gives the exact build, providers, replay instructions, and limitations.
